\setcounter{chapter}{5} % Mettre le compteur avant tout
\setcounter{section}{0}
\setcounter{figure}{0}
\setcounter{table}{0}

\counterwithin{figure}{chapter}
\counterwithin{table}{chapter}

\thispagestyle{empty}

\begin{center}
    \vspace*{\fill}

    {\fontsize{40pt}{48pt}\selectfont
    \textcolor{black}{\textsc{Chapitre 5}}}\\[2em]

    {\fontsize{22pt}{26pt}\selectfont
    \textbf{Release 3 : Cloud, Conteneurisation et DevOps}}

    \vspace*{\fill}
\end{center}

\newpage

\addcontentsline{toc}{chapter}
{Chapitre 5 : Release 3 : Cloud, Conteneurisation \& DevOps}

\fancyhead[LO,LE]
{\textbf{Chapitre 5 : Release 3 : Cloud, Conteneurisation \& DevOps}}

%====================================================
\section{Introduction}
%====================================================

Cette release a pour objectif de moderniser l’infrastructure applicative à travers l’intégration des technologies Cloud, de la conteneurisation et des pratiques DevOps.

L’objectif principal est de garantir :

\begin{itemize}[label=--]
    \item La scalabilité de l’application 
    \item La haute disponibilité des services 
    \item L’automatisation des déploiements 
    \item La portabilité des environnements 
    \item La simplification de l’administration système
\end{itemize}

Pour atteindre ces objectifs, un VPS Linux a été provisionné et configuré, les services ont été conteneurisés via Docker et orchestrés avec Docker Compose, et un pipeline CI/CD a été mis en place avec GitHub Actions.

%====================================================
\section{Sprint 6 : Déploiement sur VPS}
%====================================================

%----------------------------------------------------
\subsection{Objectif du Sprint}
%----------------------------------------------------

Ce sprint a pour objectif de mettre en place une infrastructure de production sécurisée basée sur un \textbf{VPS Linux (Ubuntu 22.04)} afin d’héberger les services applicatifs de LexAI et d’assurer une communication fiable entre les différents composants du système.

Les travaux réalisés concernent :

\begin{itemize}[label=--]
    \item Le provisionnement et la configuration initiale du VPS
    \item La configuration réseau et des règles de sécurité
    \item L’installation de l’environnement Docker sur le serveur
    \item Le déploiement des services applicatifs via Docker Compose
\end{itemize}

%----------------------------------------------------
\subsection{Présentation du VPS}
%----------------------------------------------------

Un \textbf{VPS (Virtual Private Server)} est un serveur virtuel dédié hébergé chez un fournisseur Cloud, offrant un environnement Linux isolé avec des ressources garanties.

Cette solution offre plusieurs avantages adaptés au projet LexAI :

\begin{itemize}[label=--]
    \item Environnement Linux dédié et isolé
    \item Contrôle total sur la configuration système
    \item Déploiement simplifié via SSH et Docker
    \item Coût maîtrisé pour une infrastructure de production
    \item Adresse IP publique fixe pour l’accès aux services
\end{itemize}

Dans notre projet, le VPS a été utilisé pour :

\begin{itemize}[label=--]
    \item Héberger l’ensemble des services Docker (backend, frontend, PostgreSQL, Redis, Celery)
    \item Exécuter le script de déploiement automatique (\texttt{auto-deploy.sh})
    \item Exposer l’interface LexAI via Nginx sur les ports HTTP/HTTPS
    \item Assurer la persistance des données via des volumes Docker
\end{itemize}

\begin{figure}[H]
    \centering
    \includegraphics[width=0.95\textwidth]{images/chap5 1.png}
    \caption{Vue d’ensemble de l’infrastructure VPS de production}
    \label{fig:release3-cap1}
\end{figure}

%----------------------------------------------------
\subsection{Provisionnement du VPS}
%----------------------------------------------------

Le VPS a été provisionné avec les ressources nécessaires à l’hébergement de la plateforme LexAI.

La configuration retenue comprend :

\begin{itemize}[label=--]
    \item Système d’exploitation : Ubuntu 22.04 LTS
    \item Docker Engine et Docker Compose installés
    \item Git installé pour la récupération du code source
    \item Ports ouverts : 22 (SSH), 80 (HTTP), 443 (HTTPS), 8000 (API)
\end{itemize}

Les services hébergés sur le VPS sont :

\begin{itemize}[label=--]
    \item Le backend FastAPI (conteneur \texttt{lexai-backend})
    \item Le frontend React/Nginx (conteneur \texttt{lexai-frontend})
    \item La base de données PostgreSQL (conteneur \texttt{lexai-postgres})
    \item Le broker Redis + worker Celery (conteneurs \texttt{lexai-redis}, \texttt{lexai-worker})
\end{itemize}

\begin{figure}[H]
    \centering
    \includegraphics[width=\textwidth]{images/chap5 2.png}
    \caption{Provisionnement de l’infrastructure VPS}
    \label{fig:release3-cap2}
\end{figure}

%----------------------------------------------------
\subsection{Configuration Réseau et Sécurité}
%----------------------------------------------------

Le réseau du VPS a été configuré afin d’assurer :

\begin{itemize}[label=--]
    \item L’isolation des conteneurs via le réseau Docker interne \texttt{lexai\_net}
    \item La communication sécurisée entre les services
    \item Le contrôle des flux réseau entrants et sortants
\end{itemize}

L’architecture réseau comprend :

\begin{itemize}[label=--]
    \item Un réseau Docker bridge interne (\texttt{lexai\_net}) pour la communication inter-services
    \item Un reverse proxy Nginx exposant uniquement les ports 80 et 8000
    \item Un pare-feu système bloquant tous les ports non autorisés
    \item Une adresse IP publique fixe pour l’accès externe
\end{itemize}

\begin{figure}[H]
    \centering
    \includegraphics[width=\textwidth]{images/chap5 3.png}
    \caption{Configuration réseau du VPS}
    \label{fig:release3-cap3}
\end{figure}

%----------------------------------------------------
\subsection{Règles de Sécurité}
%----------------------------------------------------

Des règles de sécurité ont été configurées afin de contrôler les accès aux services hébergés sur le VPS.

Les règles appliquées permettent :

\begin{itemize}[label=--]
    \item L’accès SSH sécurisé (port 22) pour l’administration
    \item L’accès HTTP/HTTPS (ports 80/443) pour l’interface utilisateur
    \item L’accès à l’API backend (port 8000) pour les requêtes frontend
    \item Le blocage de tous les ports non autorisés
\end{itemize}

\begin{figure}[H]
    \centering
    \includegraphics[width=\textwidth]{images/chap5 4.png}
    \caption{Règles de sécurité du VPS}
    \label{fig:release3-cap4}
\end{figure}

%----------------------------------------------------
\subsection{Déploiement des Services Applicatifs}
%----------------------------------------------------

Les services applicatifs ont été déployés sur le VPS à l’aide de Docker Compose. Le projet a été récupéré depuis GitHub, les variables d’environnement configurées dans un fichier \texttt{.env.prod}, puis les conteneurs lancés.

Les principales étapes réalisées sont :

\begin{itemize}[label=--]
    \item Connexion SSH au VPS
    \item Clonage du dépôt GitHub (\texttt{git clone})
    \item Configuration du fichier \texttt{.env.prod} (clés API, mots de passe)
    \item Lancement des services : \texttt{docker compose -f docker-compose.prod.yml up -d}
    \item Vérification de l’état des conteneurs : \texttt{docker compose ps}
\end{itemize}

\begin{figure}[H]
    \centering
    \includegraphics[width=\textwidth]{images/chap5 6.png}
    \caption{Déploiement des services applicatifs sur le VPS}
    \label{fig:release3-cap5}
\end{figure}

%----------------------------------------------------
\subsection{Clonage du Projet depuis GitHub}
%----------------------------------------------------

Le code source de LexAI a été récupéré depuis le dépôt GitHub directement sur le VPS.

\begin{figure}[H]
    \centering
    \includegraphics[width=\textwidth]{images/chap5 6.png}
    \caption{Clonage du projet LexAI depuis GitHub sur le VPS}
    \label{fig:release3-clonage}
\end{figure}

%----------------------------------------------------
\subsection{Déploiement du Backend}
%----------------------------------------------------

Le backend FastAPI a été déployé en tant que conteneur Docker sur le VPS. Les variables d’environnement (connexion PostgreSQL, Redis, clés API Groq) sont injectées via le fichier \texttt{.env.prod}.

\begin{figure}[H]
    \centering
    \includegraphics[width=\textwidth]{images/chap5 7.png}
    \caption{Déploiement du backend LexAI sur le VPS}
    \label{fig:release3-cap7}
\end{figure}

%----------------------------------------------------
\subsection{Validation de l’Infrastructure}
%----------------------------------------------------

Plusieurs tests ont été réalisés afin de vérifier le bon fonctionnement de l’infrastructure déployée :

\begin{itemize}[label=--]
    \item Accessibilité de l’interface LexAI via l’adresse IP publique du VPS
    \item Connectivité réseau entre les conteneurs Docker
    \item Communication entre le frontend, le backend, PostgreSQL et Redis
    \item Bon fonctionnement des règles de pare-feu
\end{itemize}

\begin{figure}[H]
    \centering
    \includegraphics[width=\textwidth]{images/chap5 8.png}
    \caption{Validation de la connectivité réseau et de l’accès à LexAI}
\end{figure}

%----------------------------------------------------
\subsection{Gestion du Code Source avec Git}
%----------------------------------------------------

Git a été utilisé pour gérer les versions du projet et faciliter les mises à jour du serveur. Le dépôt GitHub centralise :

\begin{itemize}[label=--]
    \item Le code source frontend et backend
    \item Les fichiers Docker et Docker Compose
    \item Le workflow GitHub Actions (\texttt{.github/workflows/deploy.yml})
    \item La documentation technique
\end{itemize}

\begin{figure}[H]
    \centering
    \includegraphics[width=\textwidth]{images/chap5 9.png}
    \caption{Gestion du code source LexAI avec Git et GitHub}
    \label{fig:release3-cap8}
\end{figure}


\section{Sprint 7 : Conteneurisation avec Docker Compose}
\subsection{Objectif du Sprint}
Ce sprint a pour objectif de conteneuriser l’ensemble des services applicatifs de LexAI à l’aide de Docker et de les orchestrer via Docker Compose afin de garantir la portabilité et la reproductibilité des déploiements.

Les travaux réalisés concernent :

\begin{itemize}[label=--]
    \item La rédaction des Dockerfiles pour chaque service
    \item La configuration du fichier \texttt{docker-compose.prod.yml}
    \item La mise en place des volumes persistants (PostgreSQL, uploads)
    \item La configuration du réseau Docker interne \texttt{lexai\_net}
\end{itemize}
\subsection{Docker et Conteneurisation}
Les services applicatifs ont été conteneurisés à l’aide de Docker afin de garantir :

\begin{itemize}[label=--]
    \item La portabilité
    \item L’isolation
    \item La standardisation des environnements
\end{itemize}

Les services conteneurisés comprennent :

\begin{itemize}[label=--]
    \item Le frontend
    \item Le backend
    \item Les services AI
    \item Certains outils techniques
\end{itemize}

\subsection{Dockerfile des Services}
Chaque service possède un Dockerfile permettant d’automatiser la génération des images applicatives.

Les Dockerfiles définissent :

\begin{itemize}[label=--]
    \item Les dépendances
    \item Les ports exposés
    \item Les commandes de démarrage
    \item L’environnement d’exécution
\end{itemize}

\begin{figure}[h]
    \centering
    \includegraphics[width=14cm]{images/dockerfile_frontend.png}
    \caption{Exemple de Dockerfile}
    \label{fig:dockerfile}
\end{figure}
\subsection{Docker Compose}
Docker Compose a été utilisé afin de simplifier le lancement simultané des services applicatifs.

Le fichier \texttt{docker-compose.yml} permet :

\begin{itemize}[label=--]
    \item La gestion des dépendances
    \item La configuration réseau
    \item Le déploiement multi-conteneurs
\end{itemize}

\begin{figure}[h]
    \centering
    \includegraphics[width=14cm]{images/dockercompose.png}
    \caption{Configuration Docker Compose}
    \label{fig:docker_compose}
\end{figure}
%====================================================
\section{Sprint 8 : Pipeline CI/CD avec GitHub Actions}
%====================================================

\subsection{Objectif du Sprint}

Ce sprint a pour objectif de mettre en place un pipeline d’intégration et de déploiement continus (CI/CD) afin d’automatiser la construction, les tests et la livraison des images Docker de la plateforme LexAI à chaque modification du code source.

Les pratiques DevOps adoptées garantissent :

\begin{itemize}[label=--]
    \item L’intégration continue des modifications de code
    \item La construction automatique des images Docker
    \item La publication des images vers un registre centralisé (GHCR)
    \item Le déploiement automatique sur le serveur de production sans intervention manuelle
    \item La traçabilité complète de chaque livraison via les tags SHA Git
\end{itemize}

%----------------------------------------------------
\subsection{Architecture du Pipeline CI/CD}
%----------------------------------------------------

Le pipeline CI/CD de LexAI repose sur \textbf{GitHub Actions}, un outil d’automatisation intégré nativement à GitHub. Il se déclenche automatiquement à chaque \textit{push} sur la branche \texttt{main}, assurant ainsi une livraison continue sans intervention manuelle.

\vspace{0.5em}

L’architecture du pipeline suit le schéma suivant :

\vspace{0.5em}

\begin{enumerate}
    \item Le développeur effectue un \texttt{git push} sur la branche \texttt{main}
    \item GitHub Actions déclenche automatiquement le workflow \texttt{deploy.yml}
    \item Le runner construit les images Docker du backend et du frontend
    \item Les images sont publiées sur le \textbf{GitHub Container Registry (GHCR)} avec deux tags : \texttt{:latest} et \texttt{:SHA-commit}
    \item Le serveur VPS exécute un script \texttt{auto-deploy.sh} (cron toutes les 5 minutes) qui détecte les nouveaux commits, effectue un \texttt{git pull} et relance les conteneurs via \texttt{docker compose up --build}
\end{enumerate}

%----------------------------------------------------
\subsection{Structure du Workflow GitHub Actions}
%----------------------------------------------------

Le fichier \texttt{.github/workflows/deploy.yml} définit le pipeline CI/CD. Il comporte un job principal \texttt{build-push} qui s’exécute sur un runner Ubuntu et réalise les étapes suivantes :

\begin{itemize}[label=--]
    \item \textbf{Checkout} : récupération du code source depuis le dépôt Git
    \item \textbf{Docker Buildx} : configuration du builder multi-plateforme
    \item \textbf{Authentification GHCR} : connexion au registre GitHub Container Registry
    \item \textbf{Build \& Push Backend} : construction de l’image \texttt{lexai-backend} et publication
    \item \textbf{Build \& Push Frontend} : construction de l’image \texttt{lexai-frontend} avec injection de la variable \texttt{VITE\_API\_BASE\_URL=/api}
\end{itemize}

\vspace{0.5em}

Le pipeline utilise le \textbf{cache de registre} (\texttt{type=registry}) pour réduire le temps de construction lors des modifications incrémentales, en réutilisant les couches Docker déjà construites.

%----------------------------------------------------
\subsection{Images Docker — Backend et Frontend}
%----------------------------------------------------

\subsubsection{Dockerfile Backend}

Le backend FastAPI est conteneurisé avec une image Python 3.11-slim. Le Dockerfile réalise les étapes suivantes :

\begin{itemize}[label=--]
    \item Installation des dépendances système (poppler, tesseract, libmagic)
    \item Installation des dépendances Python via \texttt{pip install -r requirements.txt}
    \item Copie du code source de l’application
    \item Exposition du port \texttt{8000}
    \item Démarrage du serveur Uvicorn avec rechargement automatique
\end{itemize}

\subsubsection{Dockerfile Frontend}

Le frontend React est conteneurisé en deux étapes (\textit{multi-stage build}) :

\begin{itemize}[label=--]
    \item \textbf{Étape build} : image Node 20-alpine, installation des dépendances npm, compilation Vite (\texttt{npm run build}) avec injection de \texttt{VITE\_API\_BASE\_URL}
    \item \textbf{Étape production} : image Nginx Alpine, copie des fichiers compilés depuis \texttt{dist/}, configuration du reverse proxy pour l’API backend et le routage SPA React
\end{itemize}

\begin{figure}[H]
    \centering
    \includegraphics[width=14cm]{images/docker.png}
    \caption{Dockerfile du service Frontend (multi-stage build)}
    \label{fig:dockerfile-frontend}
\end{figure}

%----------------------------------------------------
\subsection{Orchestration avec Docker Compose}
%----------------------------------------------------

Le fichier \texttt{docker-compose.yml} orchestre l’ensemble des services de la plateforme LexAI en un seul fichier de configuration déclaratif. Les services définis sont :

\begin{itemize}[label=--]
    \item \textbf{postgres} : base de données PostgreSQL 16, avec volume persistant et healthcheck
    \item \textbf{redis} : broker de messages Redis 7 pour la file de tâches Celery
    \item \textbf{backend} : service FastAPI exposé sur le port \texttt{8000}, dépendant de postgres et redis
    \item \textbf{celery-worker} : worker Celery pour l’exécution asynchrone des tâches d’analyse (Agents 1 à 4)
    \item \textbf{frontend} : service Nginx exposé sur le port \texttt{80}, servant l’interface React et proxifiant les requêtes \texttt{/api/} vers le backend
\end{itemize}

\vspace{0.5em}

Tous les services partagent le même réseau Docker interne \texttt{lexai-net} (bridge) pour communiquer de façon isolée. Les secrets (clés API, mots de passe) sont injectés via un fichier \texttt{.env} non versionné.

\begin{figure}[H]
    \centering
    \includegraphics[width=14cm]{images/dockercompose.png}
    \caption{Extrait du fichier docker-compose.yml — services backend, redis et postgres}
    \label{fig:docker-compose-full}
\end{figure}

%----------------------------------------------------
\subsection{Déploiement Automatique sur VPS}
%----------------------------------------------------

Le serveur de production est un VPS Linux (Ubuntu 22.04) hébergé sur un fournisseur Cloud. Le déploiement sur le VPS est entièrement automatisé via un script \texttt{auto-deploy.sh} exécuté toutes les 5 minutes par \texttt{cron}.

\vspace{0.5em}

Le script réalise les opérations suivantes :

\begin{enumerate}
    \item Exécution d’un \texttt{git fetch} pour récupérer les métadonnées du dépôt distant
    \item Comparaison du SHA local et du SHA distant : si identiques, aucune action
    \item En cas de nouveau commit : \texttt{git reset --hard origin/main} pour mettre à jour le code
    \item Relance des conteneurs : \texttt{docker compose up --build -d} pour reconstruire et redémarrer les services modifiés
    \item Nettoyage des images Docker obsolètes : \texttt{docker image prune -f}
\end{enumerate}

\vspace{0.5em}

Cette approche de déploiement continu (CD) basée sur le \textit{polling} Git élimine le besoin d’une connexion SSH entrante depuis GitHub Actions, contournant ainsi les restrictions de pare-feu du VPS.

%----------------------------------------------------
\subsection{Registre d’Images — GitHub Container Registry (GHCR)}
%----------------------------------------------------

Les images Docker construites par le pipeline CI sont publiées sur le \textbf{GitHub Container Registry (GHCR)}, un registre OCI intégré à GitHub. Chaque image est taguée avec deux identifiants :

\begin{itemize}[label=--]
    \item \texttt{:latest} — image la plus récente, utilisée pour le déploiement
    \item \texttt{:SHA-commit} — identifiant unique du commit Git, permettant un retour en arrière (\textit{rollback}) précis vers n’importe quelle version antérieure
\end{itemize}

\vspace{0.5em}

L’authentification au registre est réalisée via le token \texttt{GITHUB\_TOKEN} généré automatiquement par GitHub Actions, sans nécessiter de secret supplémentaire.

%----------------------------------------------------
\subsection{Bilan DevOps}
%----------------------------------------------------

\begin{table}[H]
\centering
\small
\renewcommand{\arraystretch}{1.3}
\begin{tabular}{|>{\RaggedRight\arraybackslash}p{3.5cm}|>{\RaggedRight\arraybackslash}p{4cm}|>{\RaggedRight\arraybackslash}p{5cm}|}
\hline
\textbf{Pratique DevOps} & \textbf{Outil utilisé} & \textbf{Bénéfice} \\
\hline
Intégration Continue (CI) & GitHub Actions & Détection automatique des erreurs de build \\
\hline
Livraison Continue (CD) & GHCR + cron VPS & Déploiement sans intervention manuelle \\
\hline
Conteneurisation & Docker + Dockerfile & Portabilité et reproductibilité \\
\hline
Orchestration locale & Docker Compose & Gestion multi-services simplifiée \\
\hline
Gestion des secrets & Fichier \texttt{.env} & Sécurisation des clés API et mots de passe \\
\hline
Versionnage images & Tag SHA Git & Rollback précis vers toute version \\
\hline
Cache build & Docker layer cache & Réduction du temps de build CI \\
\hline
\end{tabular}
\caption{Bilan des pratiques DevOps adoptées pour LexAI}
\label{tab:devops-bilan}
\end{table}

\newpage

%====================================================
\subsection{Conclusion}
Cette release a permis de mettre en place une infrastructure de production moderne basée sur un VPS Linux, Docker et un pipeline CI/CD automatisé avec GitHub Actions.

L’utilisation de ces technologies a permis :

\begin{itemize}[label=--]
    \item D’améliorer la portabilité des services grâce à la conteneurisation Docker
    \item De standardiser les environnements de développement et de production
    \item D’assurer une meilleure scalabilité de l’infrastructure via Docker Compose
    \item De simplifier et d’automatiser le déploiement des applications grâce au pipeline CI/CD
    \item De garantir la traçabilité complète de chaque livraison via les tags SHA Git
\end{itemize}

L’architecture DevOps mise en œuvre constitue une base robuste et évolutive, adaptée à la croissance future de la plateforme LexAI.